

Os requisitos são especificações do sistema que visam suprir necessidades de usuários, sendo a elicitação de requisitos uma etapa crucial para o desenvolvimento de software. Esta é uma etapa fundamental para a engenharia de requisitos \cite{reinehr2020engenharia}. Segundo \citeonline{sommerville_2011}, esse processo deve ser realizado de forma sistemática com diversas técnicas, como documentação, validação, negociação, entre outras.

Sendo assim, para o projeto proposto foram realizadas personas, brainstorming e histórias de usuário para elicitar novos requisitos. Além disso, MoSCoW e revisão de funcionalidades foram utilizadas para priorização dos requisitos levantados e a classificação dos requisitos obtidos foi realizada utilizando o método FURPS+, proporcionando uma visão ampla das necessidades dos usuários e do sistema.

\subsection{Requisitos Funcionais}

Os requisitos funcionais representam as funcionalidades que o sistema deve oferecer para atender às necessidades explícitas dos usuários. Em essência, eles descrevem comportamentos e serviços que o sistema deverá executar, além de regras de negócio relacionadas ao contexto do sistema, e sua implementação resulta diretamente em entregas perceptíveis ao usuário final \cite{reinehr2020engenharia}.

\subsubsection{Histórias de Usuário}

As histórias de usuário têm como principal objetivo descrever, de forma clara e concisa, funcionalidades necessárias do ponto de vista do usuário. Outro aspecto importante é a possibilidade de dividi-las em épicos, sendo grandes conjuntos de funcionalidades relacionadas que representam macro-requisitos do sistema, e \textit{features}, sendo funcionalidades específicas derivadas dos épicos, com escopo mais restrito e bem definido \cite{reinehr2020engenharia}.

Além disso, a fim de proporcionar maior controle sobre as funcionalidades criadas, é necessário um método de priorização. O método escolhido para este projeto foi o MoSCoW, que classifica os requisitos funcionais nas seguintes categorias \cite{carvalho2021moscow}:

\begin{itemize}
    \item \textbf{Must}: requisitos que o sistema obrigatoriamente deve ter;
    \item \textbf{Should}: requisitos importantes, mas que podem ser postergados caso necessário;
    \item \textbf{Could}: requisitos desejáveis, cuja implementação depende da disponibilidade de tempo e recursos;
    \item \textbf{Won’t}: requisitos que, embora desejados, não serão implementados neste momento.
\end{itemize}

Sendo assim, os tópicos a seguir representam funcionalidades desenvolvidas usando histórias de usuário e priorizadas com MoSCoW. Cada funcionalidade é identificada por um código no formato \textbf{FX} para representar uma \textit{feature}, e \textbf{USX} para indicar a respectiva história de usuário associada.

\subsubsubsection{Épico - Gerenciamento de registro do usuário.}

\begin{itemize}

    \item \textbf{Feature F01}: registro do usuário.
    
    US1 Eu, como usuário, desejo me cadastrar no aplicativo.

    \begin{itemize}
        \item Prioridade: deve ter.
    \end{itemize}
    
    \item \textbf{Feature F02}: registro do funcionário.

    US2 Eu, como funcionário, desejo me cadastrar no aplicativo.

    \begin{itemize}
        \item Prioridade: deve ter.
    \end{itemize}
    
    \item \textbf{Feature F03}: edição de dados do usuário ou funcionário.

    US3 Eu, como usuário ou funcionário, desejo editar meus dados no aplicativo.

    \begin{itemize}
        \item Prioridade: deve ter.
    \end{itemize}

    \item \textbf{Feature F04}: exclusão de dados do usuário ou funcionário.

    US4 Eu, como usuário ou funcionário, desejo deletar os meus dados do aplicativo.

    \begin{itemize}
        \item Prioridade: deve ter.
    \end{itemize}
    
\end{itemize}

%--------------------------------------------------------------------------

\subsubsubsection{Épico - Gerenciamento de formulários.}

\begin{itemize}

    \item \textbf{Feature F05}: registro de formulário.

    US5 Eu, como usuário, desejo registrar diariamente minha avaliação do meu estado em cada seção do bem-estar.

    \begin{itemize}
        \item Prioridade: deve ter.
    \end{itemize}
    
    US6 Eu, como funcionário, desejo visualizar os dados registrados pelos usuários.

    \begin{itemize}
        \item Prioridade: deve ter.
    \end{itemize}
    
    \item \textbf{Feature F06}: visualização de dados do formulário.

    US7 Eu, como usuário, desejo visualizar meus formulários respondidos a qualquer momento.

    \begin{itemize}
        \item Prioridade: deveria ter.
    \end{itemize}
    
    \item \textbf{Feature F07}: estatísticas de formulários respondidos.

    US8 Eu, como usuário, desejo visualizar as estatísticas dos formulários respondidos em forma de gráfico a qualquer momento.

    \begin{itemize}
        \item Prioridade: deveria ter.
    \end{itemize}
    
    \item \textbf{Feature F08}: gerenciamento de formulários.

    US9 Eu, como funcionário, desejo criar formulários personalizados.

    \begin{itemize}
        \item Prioridade: poderia ter.
    \end{itemize}

    US10 Eu, como funcionário, desejo editar formulários à medida que analiso as perguntas.

    \begin{itemize}
        \item Prioridade: deveria ter.
    \end{itemize}

    \item \textbf{Feature F09}: bloco de anotações.

    US11 Eu, como usuário, desejo adicionar anotações livres para os formulários respondidos.

    \begin{itemize}
        \item Prioridade: poderia ter.
    \end{itemize}

\end{itemize}

%--------------------------------------------------------------------------

\subsubsubsection{Épico - Gerenciamento de comunicação.}

\begin{itemize}
    \item \textbf{Feature F010}: geração de notificações.

    US12 Eu, como usuário, gostaria de receber notificações para preenchimento dos formulários diários.

    \begin{itemize}
        \item Prioridade: poderia ter.
    \end{itemize}
    
    \item \textbf{Feature F11}: troca de senha.

    US13 Eu, como usuário, desejo receber a confirmação de troca de senha por e-mail.

    \begin{itemize}
        \item Prioridade: deveria ter.
    \end{itemize}
    
    \item \textbf{Feature F12}: gerenciamento de avisos.

    US14 Eu, como usuário, gostaria de receber avisos sobre o status do aplicativo por e-mail.
    
    \begin{itemize}
        \item Prioridade: poderia ter.
    \end{itemize}
    
    \item \textbf{Feature F13}: avisos de eventos.

    US15 Eu, como usuário, gostaria de receber avisos sobre possíveis eventos dentro do calendário do aplicativo.
    
    \begin{itemize}
        \item Prioridade: poderia ter.
    \end{itemize}
\end{itemize}

%--------------------------------------------------------------------------

\subsubsubsection{Épico - Gerenciamento de eventos.}

\begin{itemize}
    \item \textbf{Feature F014}: gerenciar evento.

    US17 Eu, como funcionário, desejo ser capaz de registrar eventos no calendário do aplicativo.

    \begin{itemize}
        \item Prioridade: deveria ter.
    \end{itemize}
    
    US18 Eu, como funcionário, desejo ser capaz de editar informações de eventos no calendário do aplicativo.

    \begin{itemize}
        \item Prioridade: deveria ter.
    \end{itemize}
    
    US19 Eu, como funcionário, desejo ser capaz de excluir eventos no calendário do aplicativo.

    \begin{itemize}
        \item Prioridade: deveria ter.
    \end{itemize}
    
    \item \textbf{Feature F15}: visualizar evento.

    US20 Eu, como usuário, desejo visualizar eventos no calendário do meu aplicativo.

    \begin{itemize}
        \item Prioridade: deveria ter.
    \end{itemize}

\end{itemize}

%--------------------------------------------------------------------------

\subsubsubsection{Épico - Gerenciamento do feed.}

\begin{itemize}
    \item \textbf{Feature F016}: gerenciar feed.

    US21 Eu, como funcionário, desejo ser capaz de registrar publicações no feed de cada seção do bem-estar.

    \begin{itemize}
        \item Prioridade: deveria ter.
    \end{itemize}

    US22 Eu, como funcionário, desejo ser capaz de editar publicações no feed de cada seção do bem-estar.

    \begin{itemize}
        \item Prioridade: deveria ter.
    \end{itemize}

    US23 Eu, como funcionário, desejo ser capaz de excluir publicações no feed de cada seção do bem-estar.

    \begin{itemize}
        \item Prioridade: deveria ter.
    \end{itemize}

    \item \textbf{Feature F017}: visualizar o feed.

    US24 Eu, como usuário, desejo ser capaz de visualizar o feed para cada seção do bem-estar.
    
    \begin{itemize}
        \item Prioridade: deveria ter.
    \end{itemize}

\end{itemize}

Sendo assim, com base nos requisitos elicitados pelo Lean Inception com Personas e Brainstorming de Funcionalidades, juntamente com os requisitos elicitados pelas Histórias de Usuário, priorizados por Revisão de Funcionalidades e MoSCoW, logo a tabela ~\ref{tab:requisitosfunc} representa os requisitos funcionais para o projeto proposto.


\begin{center}
\begin{longtable}{|>{\raggedright\arraybackslash}p{2cm}|p{13cm}|}
\caption{Requisitos Funcionais} \label{tab:requisitosfunc} \\
\hline
\textbf{ID} & \textbf{Descrição} \\
\hline
\endfirsthead

\hline
\textbf{ID} & \textbf{Descrição} \\
\hline
\endhead

        FUNC01 & Cadastro de usuário no aplicativo. \\
        \hline
        FUNC02 & Cadastro de funcionário no aplicativo. \\
        \hline
        FUNC03 & Edição dos dados cadastrais de usuários e funcionários. \\
        \hline
        FUNC04 & Exclusão dos dados de usuários e funcionários. \\
        \hline
        FUNC05 & Registro diário de formulários pelos usuários. \\
        \hline
        FUNC06 & Visualização de dados registrados por usuários (funcionário). \\
        \hline
        FUNC07 & Visualização dos formulários respondidos (usuário). \\
        \hline
        FUNC08 & Geração de estatísticas dos formulários respondidos. \\
        \hline
        FUNC09 & Criação e edição de formulários personalizados (funcionário). \\
        \hline
        FUNC10 & Inclusão de anotações livres associadas aos formulários. \\
        \hline
        FUNC11 & Envio de notificações para preenchimento dos formulários. \\
        \hline
        FUNC12 & Confirmação de troca de senha via e-mail. \\
        \hline
        FUNC13 & Envio de avisos sobre o status do aplicativo por e-mail. \\
        \hline
        FUNC14 & Envio de avisos sobre eventos no calendário do aplicativo. \\
        \hline
        FUNC15 & Registro de eventos no calendário (funcionário). \\
        \hline
        FUNC16 & Edição e exclusão de eventos do calendário (funcionário). \\
        \hline
        FUNC17 & Visualização dos eventos no calendário (usuário). \\
        \hline
        FUNC18 & Registro de publicações no feed por seção do bem-estar. \\
        \hline
        FUNC19 & Edição e exclusão de publicações no feed (funcionário). \\
        \hline
        FUNC20 & Visualização do feed por seção do bem-estar (usuário). \\
        \hline
        FUNC21 & Exportação dos dados dos formulários via CSV. \\
        \hline
        FUNC22 & Login com Google (OAuth 2.0). \\
        \hline
        FUNC23 & Geração de análises e estatísticas instituicionais. \\
        \hline
        
\multicolumn{2}{r}{\small Fonte: autoria própria.} \\
\end{longtable}
\end{center}


%--------------------------------------------------------------------------

\subsection{Requisitos Não Funcionais.}

Requisitos não funcionais são aqueles que não são mapeados facilmente por meios convencionais, ou seja, são características do sistema que não estão relacionadas diretamente com as necessidades dos usuários \cite{sommerville_2011}. Sendo assim, \citeonline{sommerville_2011} classifica esses requisitos como desempenho, segurança, usabilidade, disponibilidade, entre outros, sendo esses não perceptíveis pelo usuário, mas essenciais para o sistema.

\subsubsection{FURPS+}

FURPS é um modelo de classificação de requisitos de software criado para o Rational Unified Process, cujo principal objetivo é o conjunto de melhores práticas para organizar o processo de construção de software. Sendo assim, no FURPS, os requisitos são divididos em cinco principais categorias: (F)unctionality, (U)sability, (R)eliability, (P)erformance, (S)upportability \cite{campos_2008}.

Além disso, existe uma extensão conhecida como \textbf{FURPS+}, onde é implementado o FURPS convencional com um tipo extra de classificação de requisitos, podendo ser requisitos de interfaces, implementações e físicos, como restrições de hardware, por exemplo \cite{sequeira2017usabilidade}.


\begin{itemize}

    \item \textbf{F: Funcionalidade}: Aspectos funcionais do software, esses requisitos foram mapeados pelas histórias de usuário, personas, entre outros. 

    \item \textbf{U: Usabilidade}: Usabilidade do sistema, incluindo experiência do usuário (UX).

    \item \textbf{R: Confiabilidade}: Confiabilidade do sistema, como estabilidade e resistência a falhas.

    \item \textbf{P: Performance}: Performance do sistema, como velocidade e eficiência.

    \item \textbf{S: Suportabilidade}: Avalia a facilidade com que o sistema pode ser mantido e suportado ao longo do tempo.

    \item \textbf{+}: Requisitos como ajuda e documentação de usuário, interfaces, observações legais, entre outros.

\end{itemize}

Sendo assim, a partir das especificações de requisitos não funcionais propostas por \citeonline{sommerville_2011} e necessidades do sistema, é possível identificar os seguintes requisitos não funcionais para o projeto proposto:

\renewcommand{\arraystretch}{1.5}
\begin{table}[H]
    \centering
    \caption{Requisitos Não Funcionais Classificados Pelo FURPS+.}
    \label{tab:requisitosnaofunc}
    
    \begin{tabularx}{\textwidth}{|c|X|}
        \hline
        \textbf{ID} & \textbf{Descrição} \\
        \hline
        USAB01 & O sistema deve oferecer suporte a vários idiomas.  \\
        \hline
        USAB02 & O app deve responder imediatamente às ações do usuário, mesmo com carregamentos. \\
        \hline
        USAB03 & O sistema deve seguir um padrão simples para facilitar o aprendizado. \\
        \hline
        USAB04 & As interações devem ter feedback visual ou sonoro para orientar o usuário. \\
        \hline
        USAB05 & O app deve mostrar mensagens claras em caso de erro, sem atrapalhar o uso. \\
        \hline
        CONF01 & O usuário pode alterar seus dados pessoais. \\
        \hline
        CONF02 & O app deve seguir a Lei Geral de Proteção de Dados (LGPD). \\
        \hline
        CONF03 & O sistema deve funcionar 24 horas por dia, 7 dias por semana. \\
        \hline
        CONF04 & Deve existir um plano de backup para recuperação em caso de falhas. \\
        \hline
        PERF01 & O sistema deve suportar acessos simultâneos de diferentes usuários. \\
        \hline
        PERF02 & O aplicativo deve realizar autenticação no tempo médio de 10 segundos. \\
        \hline
        PERF03 & O tempo de resposta deve ser inferior a 2 segundos. \\
        \hline
        PERF04 & O sistema deve aguentar um aumento de 20\% no número de usuários. \\
        \hline
        SUPO01 & O app deve estar disponível para Android e iOS, com interface semelhante. \\
        \hline
        SUPO02 & O app deve se adaptar a diferentes dispositivos, como celulares e tablets. \\
        \hline
        SUPO03 & O sistema deve seguir padrões de segurança reconhecidos, como ISO 27001. \\
        \hline
        PLUS01 & O sistema deve oferecer suporte online, dicas e manual acessível. \\
        \hline
    \end{tabularx}
    
    \vspace{0.5em}
    \small Fonte: autoria própria.
\end{table}



